\chapter{Currency−Free Economy: XP, Barter, and Trusts}
łabel{ch:currency−free−economy}

\begin{tcolorbox}[colback=blue!5!white,colframe=blue!70!black,title=The Ledger of Memory]
\emph{``The world does not run on silver. It runs on memory−−−on favors owed, secrets kept, and the weight of promises. A coin can be stolen; a debt must be paid in deeds.''}  
−−− Dusana of the Raven Road, \textit{The Hearth Ledger}
\end{tcolorbox}

\section{No Coin, Only Consequence}
łabel{sec:no−coin}

In Fate's Edge, there is no abstract currency. No gold pieces, no silver scales, no ledgers of mundane wealth. What you \emph{know}, whom you \emph{know}, and what you \emph{owe} matter more than what you carry.

\begin{itemize}
  ℹtem \textbf{Incidental expenses} −− a meal, a night at an inn, a riding horse, a week's trail rations −− are \textbf{hand‑waved}. Assume your character has enough from past earnings, barter, or simple hospitality to cover ordinary life. If the fiction demands scarcity, the GM will introduce a Supply Timer or a Complication.
  ℹtem \textbf{Significant acquisitions} −− assets, favors, magical items, political influence, or the service of a powerful ally −− require one of three forms of payment:
  \begin{enumerate}
    ℹtem \textbf{Experience Points (XP)} −− Spend your hard‑won growth to secure lasting resources.
    ℹtem \textbf{Barter or Strings} −− Trade goods, services, information, or narrative leverage (Strings) for equivalent value.
    ℹtem \textbf{Obligation} −− Accept a Debt Timer or increase Obligation to a patron (terrestrial or supernatural) in exchange for immediate gain.
  \end{enumerate}
\end{itemize}

This chapter explains how to spend XP on non‑personal advancement, how to pool XP for party assets, how to use Strings and Obligation as currency, and how to form a player‑owned terrestrial patron: a \textbf{Trust}.

\section{Abstracted Treasure: Finds That Change the Story}
łabel{sec:abstracted−treasure}

You do not track copper or silver. Instead, treasure is \textbf{leverage} −−− each meaningful find has a one‑time narrative effect. A bag of gold is not treasure. A gold coin that was once paid as a traitor's blood‑price \emph{is}.

\subsection{Treasure Tiers}
łabel{subsec:treasure−tiers}

\begin{center}
\begin{tabular}{lll}
⊤rule
\textbf{Tier} & \textbf{Example} & \textbf{What It Does} \\
∣rule
\textbf{Minor} & A merchant's seal, a noble's letter, a pouch of star‑iron beads & Buy a single favor: safe passage for one night, a small bribe, a rumor, ignore a minor toll \\
\textbf{Standard} & A signed treaty, a bridge's deed, a saint's relic, a ship's luck in a bottle & Buy a significant exception: clear a crime, force a parley, open a sealed gate, gain an audience with a lord \\
\textbf{Major} & A king's true name, a lullaby that calms armies, a dragon's scale, a prophecy sealed in amber & Change the story: rewrite a local law, end a feud, call a truce, redirect a disaster \\
⊥tomrule
\end{tabular}
\end{center}

\subsection{Finding and Using Treasure}
łabel{subsec:using−treasure}

When players loot a significant location (a dragon's lair, a fallen noble's vault, a sunken temple), the GM awards \textbf{1–3 treasure items} at appropriate tiers. No coin. No weight. No bookkeeping.

A character may carry at most \textbf{3 treasure items} total (any mix). Unused treasure fades after one session unless invested in an Asset (pay the treasure's XP equivalent to keep it as a permanent String).

Spend a treasure item to:

\begin{itemize}
  ℹtem \textbf{Improve Position} on a relevant roll by one step (Minor).
  ℹtem \textbf{Auto‑succeed} on a non‑combat roll where leverage applies (Standard).
  ℹtem \textbf{Rewrite a scene's stakes} (Major) −−− e.g., ``I use the king's true name to make his guards stand down.''
\end{itemize}

\subsection{Treasure as XP (Optional)}
łabel{subsec:treasure−as−xp}

A player may convert a treasure item into XP instead of using its narrative effect:

\begin{itemize}
  ℹtem Minor $→$ 1 XP
  ℹtem Standard $→$ 2 XP
  ℹtem Major $→$ 4 XP
\end{itemize}

The GM should warn: \emph{``You can sell the dragon's scale for power, or you can use it to command the dragon's respect.''}

\subsection{Integration with Existing Systems}
łabel{subsec:treasure−integration}

\begin{itemize}
  ℹtem \textbf{Assets:} A treasure item can be spent to acquire a Minor Asset instantly (no downtime).
  ℹtem \textbf{Debt Timers:} A treasure item can clear 1–3 segments of a Debt Timer (GM discretion).
  ℹtem \textbf{Strings:} A treasure item can become a permanent String if invested during downtime.
  ℹtem \textbf{Boons:} A treasure item can be cashed for 2 Boons instead of its narrative effect.
\end{itemize}

\paragraph{The Golden Rule of Treasure:}
\emph{If it doesn't change the story, it's not treasure.} Name it. Describe it. Then watch your players hoard it not for its value, but for its power to change what happens next.

\begin{tcolorbox}[colback=yellow!5!white,colframe=yellow!70!black,title=Abstracted Treasure: The Leverage Standard]
  \emph{``A Minor treasure affects a single roll or scene. A Standard treasure affects a single conflict or session. A Major treasure changes the campaign arc or permanently rewrites the fiction.''}
  
  \vspace{2mm}
  
  \noindent
  \begin{tabular}{p{2.5cm}p{5cm}p{4.5cm}}
  ⊤rule
  \textbf{Tier} & \textbf{Fate's Edge Example} & \textbf{What It Does} \\
  ∣rule
  μltirow{3}{2.5cm}{Minor} 
  & \textit{Emberfang Dagger} & Glows when undead are near; once per scene, gain +1 die to detect hidden spirits. \\
  & \textit{Vial of Starlight Water} & Drink to clear 1 Fatigue or ignore a minor poison (one use). \\
  & \textit{Runes of Unsteady Passage} & Chalk a doorframe; the next person through forgets why they entered (one scene). \\
  ∣rule
  μltirow{3}{2.5cm}{Standard} 
  & \textit{Bell of the Fog Warden} & Ring it to part mist or silence a supernatural wail for the entire scene. \\
  & \textit{Captain's Luck in a Bottle} & Break it to reroll any failed roll for your ship or crew (once per session). \\
  & \textit{Sealed Knight's Confession} & Read it aloud to force a parley; the target cannot attack for the duration of the scene. \\
  ∣rule
  μltirow{3}{2.5cm}{Major} 
  & \textit{True Name of the Ashen Tyrant} & Speak it to compel the tyrant to abdicate, or to force open any gate they sealed. \\
  & \textit{Seed of the First Oak} & Bury it to grow a fortress in a day −−− the fortress remembers the builder's fears. \\
  & \textit{Fossilised Heart of the Elder Wyrm} & Once per campaign, negate a natural disaster or redirect a region‑wide curse. \\
  ⊥tomrule
  \end{tabular}
  
  \vspace{2mm}
  \noindent
  \textbf{Leverage Test:} Ask \emph{``What can this treasure do \textbf{once}?''} \\
    −−− ``Give +1 die or skip a minor obstacle'' $→$ Minor \\
    −−− ``Resolve a significant problem or open a major opportunity'' $→$ Standard \\
    −−− ``End a war, topple a dynasty, or redirect destiny'' $→$ Major
  
  \vspace{1mm}
  \noindent
  \textit{Passive, permanent benefits} (e.g., a sword that always adds +1 die) are \textbf{Assets}, not treasures. The treasure version is a one‑scene surge of that power.
\end{tcolorbox}

\section{What Costs XP (and What Doesn't)}
łabel{sec:xp−costs−table}

\subsection{Hand‑Waved (Free)}
łabel{subsec:handwaved}

Always assume you have enough for:
\begin{itemize}
  ℹtem Ordinary meals, lodging, basic travel rations.
  ℹtem A riding horse or pack mule (if fictionally plausible).
  ℹtem Basic tools (rope, lantern, writing kit, lockpicks).
  ℹtem Minor repairs, small bribes (a drink, a gift).
  ℹtem Routine transportation (ferry, caravan passage) \emph{unless} the journey itself is the adventure.
\end{itemize}

\paragraph{GM Guideline:} When a player asks ``Do I have enough for X?'' and X is not a major story resource, answer ``Yes'' unless scarcity is the point of the scene.

\subsection{Costs XP (or Equivalent Barter / Obligation)}
łabel{subsec:xp−cost−items}

\begin{center}
\begin{tabular}{lll}
⊤rule
\textbf{Item / Benefit} & \textbf{XP Cost (min.)} & \textbf{Notes} \\
∣rule
Minor Asset (safehouse, workshop, contact network) & 4 XP & See \ref{sec:asset−types−costs} \\
Standard Asset (guild seat, spy ring, fortress lease) & 8 XP & See \ref{sec:asset−types−costs} \\
Major Asset (chartered city license, fleet, noble title) & 12 XP & See \ref{sec:asset−types−costs} \\
Follower (Capability squared) & Cap² XP & See \ref{sec:follower−costs} \\
Magical item (single use) & 2−−4 XP & GM discretion \\
Magical item (permanent, Tier I) & 6−−10 XP & Requires downtime \\
Clearing a Debt Timer segment & 2 XP per segment & Instead of performing a service \\
Buying a permanent String (favor owed to you) & 2−−4 XP & One‑time narrative permission \\
Upkeep of an Asset or Follower (efficient method) & As per \ref{subsec:upkeep−options} & XP instead of coin \\
Forming a Trust (player‑owned patron) & 10 XP (party pool) & See \ref{sec:forming−trust} \\
⊥tomrule
\end{tabular}
\end{center}

\paragraph{Design Note:} XP represents not just training but \emph{investment} −− resources, contacts, and time spent securing something of lasting value. Want a safehouse? You don't pay gold; you spend XP to reflect the effort of finding, securing, and stocking it.

\section{Pooling XP for Party Assets}
łabel{sec:party−pool}

A party may create a \textbf{shared pool} of XP to acquire and maintain assets that benefit everyone. The pool is tracked separately from individual XP.

\subsection{Creating a Pool}
łabel{subsec:creating−pool}

\begin{itemize}
  ℹtem Any player may contribute any amount of their unspent XP to the party pool.
  ℹtem Contributions are voluntary and irreversible (unless the party agrees to dissolve the pool −− see \ref{subsec:dissolving−pool}).
  ℹtem The pool has no owner; decisions about spending from it require \textbf{majority vote} (or table consensus).
\end{itemize}

\subsection{Spending from the Pool}
łabel{subsec:spending−pool}

\begin{itemize}
  ℹtem The pool can only be used for \textbf{party assets} (shared safehouse, group ship, guild charter, spy network) or for \textbf{forming a Trust} (see \ref{sec:forming−trust}).
  ℹtem It cannot be used for individual Attribute/Skill increases, personal Talents, or personal Followers.
  ℹtem When an asset bought with the pool requires \textbf{upkeep} (see \ref{subsec:upkeep−options}), the upkeep can be paid from the pool (efficient method) or by assigning a party member to handle it (intensive method).
\end{itemize}

\subsection{Dissolving the Pool}
łabel{subsec:dissolving−pool}

\begin{itemize}
  ℹtem If the party permanently disbands or agrees to liquidate, the remaining XP in the pool is \textbf{lost} −− it represents sunk investment in shared infrastructure.
  ℹtem Instead, the party may \textbf{transfer} the asset to one character (who must pay its XP cost from their own XP to ``buy out'' the others). The other contributors are compensated with an equivalent String from that character (see \ref{sec:barter−strings−obligation}).
\end{itemize}

\paragraph{Example:} Alexander (9 XP banked), Lyra (4 XP), and Sera (6 XP) want a shared safehouse (Minor Asset, 4 XP). They contribute 2 XP each to the pool (total 6 XP). They spend 4 XP to acquire the safehouse, leaving 2 XP in the pool for future upkeep. The safehouse is now a party asset. Later they add Mira (3 XP), who contributes 1 XP to the pool.

\section{Barter, Strings, and Obligation}
łabel{sec:barter−strings−obligation}

When you lack XP, you may pay with \textbf{barter} (trading goods or services), \textbf{Strings} (narrative leverage), or \textbf{Obligation} (debt to a patron). These are not coin −− they are story fuel.

\subsection{Barter}
łabel{subsec:barter}

\begin{itemize}
  ℹtem Trade an item of comparable value (a magical sword for a spy network, a noble's secret for a ship).
  ℹtem The GM adjudicates fairness. A barter exchange usually requires a Sway or Broker roll (DV 2−−4).
  ℹtem \textbf{Service Barter:} Promise a future task in exchange for immediate gain. This creates a Debt Timer [4−−6] for the recipient.
\end{itemize}

\subsection{Strings as Currency}
łabel{subsec:strings−currency}

Strings are defined in the core rules as narrative permissions −− a favor owed, a secret known, a charter held. They can be \textbf{spent} as currency:

\begin{center}
\begin{tabular}{ll}
⊤rule
\textbf{String Value} & \textbf{Can purchase} \\
∣rule
Minor String (a small favor, a local contact) & One‑use discount of 1 XP on an asset, or a single use of an Asset's off‑screen effect \\
Standard String (a noble's oath, a guild seal) & A Minor Asset (4 XP value) or clear 2 segments of a Debt Timer \\
Major String (a life‑debt, a royal charter) & A Standard Asset (8 XP value) or a Follower (up to Cap 3) \\
⊥tomrule
\end{tabular}
\end{center}

To convert a String into XP, the player must \textbf{use} the String in play −− narrate how it secures the asset. The GM may require a roll (e.g., Presence + Sway to call in the favor).

\subsection{Obligation as Payment}
łabel{subsec:obligation−payment}

Instead of XP, accept a \textbf{Debt Timer } or increase \textbf{Obligation} to a terrestrial patron (see \ref{sec:terrestrial−patrons}).

\begin{center}
\begin{tabular}{ll}
⊤rule
\textbf{Debt Timer Size} & \textbf{Equivalent Value} \\
∣rule
[4] & Minor Asset or one‑time service \\
[6] & Standard Asset or Cap 3 Follower \\
[8] & Major Asset or Cap 5 Follower \\
⊥tomrule
\end{tabular}
\end{center}

You may clear a Debt Timer segment by spending 2 XP (as if buying off the debt with hard work), or by performing an appropriate service (GM discretion).

\paragraph{Example:} The party wants a ship (Standard Asset, 8 XP). They have no XP, but a Sidhi pirate captain offers them the ship in exchange for a \textbf{Debt Timer [6]}. When it fills, they must raid an Oshiiran grain barge for her. They accept.

\section{Player‑Owned Terrestrial Patrons: Trusts}
łabel{sec:trusts}

A \textbf{Trust} is a collective entity −− a guild, a covenant, a company, a secret society −− founded and run by the player characters. It acts as a \textbf{terrestrial patron} (see \ref{sec:terrestrial−patrons}), but the players collectively decide its demands and benefits.

\subsection{Forming a Trust}
łabel{sec:forming−trust}

\begin{itemize}
  ℹtem Requires \textbf{10 XP from the party pool} (cannot be paid individually).
  ℹtem The party names the Trust, defines its purpose (e.g., ``The Silkhand League −−− spice merchants and smugglers''), and chooses an initial \textbf{Seat} (headquarters or symbol).
  ℹtem The Trust starts with a \textbf{Debt Timer [4]} (the cost of founding) that the players must clear through roleplay or XP before it can grant benefits.
\end{itemize}

\subsection{How a Trust Works}
łabel{subsec:trust−mechanics}

The Trust is a \textbf{non‑player character} entity with:

\begin{itemize}
  ℹtem \textbf{Obligation Capacity} equal to the party's combined Spirit + Presence (or a fixed value of 6 for simplicity). Members incur Obligation to the Trust when they ask for its help.
  ℹtem \textbf{Debt Timer } (size 4, 6, or 8) −− tracks external obligations the Trust owes to other factions.
  ℹtem \textbf{Asset Pool} −− the Trust can own Assets (safehouses, ships, contact networks) bought with party XP (see \ref{sec:party−pool}).
  ℹtem \textbf{Follower Pool} −− the Trust can employ Followers (e.g., guards, clerks, spies) paid from the party pool.
\end{itemize}

\subsection{Benefits of Belonging to a Trust}
łabel{subsec:trust−benefits}

A character who is a member of the Trust (all founders are members; new members can be inducted) may:

\begin{itemize}
  ℹtem \textbf{Call on the Trust} −− once per session, gain +1 die on a roll directly related to the Trust's purpose (e.g., smuggling, trade, information). This costs no Obligation.
  ℹtem \textbf{Borrow an Asset} −− use one of the Trust's Assets for a scene. Mark +1 Obligation to the Trust per Asset used.
  ℹtem \textbf{Request a Follower} −− temporarily command a Trust follower (Cap 1−−2) for a mission. Mark +1 Obligation per follower.
  ℹtem \textbf{Use the Trust's Strings} −− the Trust may accumulate Strings (favors from others). A member may spend one such String as if it were their own (costs 1 Obligation).
\end{itemize}

\subsection{Obligation to the Trust}
łabel{subsec:trust−obligation}

\begin{itemize}
  ℹtem A member's Obligation to the Trust is tracked separately from other patrons. Capacity = Spirit + Presence (see \ref{sec:obligation−capacity}).
  ℹtem When Obligation fills (exceeds capacity), the Trust \textbf{demands a task} −− a mission that advances its goals. If refused, the member is \textbf{expelled} and loses all Trust benefits.
  ℹtem Obligation can be cleared by:
  \begin{itemize}
    ℹtem Performing a service for the Trust (GM discretion, usually clears 1−−2 segments).
    ℹtem Contributing XP to the Trust pool (2 XP per segment).
    ℹtem Donating a personal Asset or String of equivalent value.
  \end{itemize}
\end{itemize}

\subsection{Growing the Trust}
łabel{subsec:trust−growth}

The party may invest additional XP from the party pool to expand the Trust:

\begin{center}
\begin{tabular}{ll}
⊤rule
\textbf{Investment} & \textbf{Effect} \\
∣rule
4 XP & Increase Debt Timer capacity by 2 (can owe more before demanding repayment) \\
6 XP & Add a new Trust Asset (Minor) \\
8 XP & Upgrade an Asset from Minor to Standard \\
10 XP & Add a Trust Follower (Cap 3) \\
12 XP & Gain a permanent String (a faction owes the Trust a favor) \\
15 XP & The Trust becomes a \textbf{Major Terrestrial Patron} −− can grant its own Rites (narrative permissions) \\
⊥tomrule
\end{tabular}
\end{center}

\subsection{Trusts as Story Engines}
łabel{subsec:trust−story}

The GM should use the Trust as a source of complications:

\begin{itemize}
  ℹtem \textbf{Schism} −− A faction within the Trust disagrees with leadership. Start a Trust Loyalty timer [6].
  ℹtem \textbf{Hostile Takeover} −− Another group (rival guild, Ashaani agents) tries to seize the Trust's assets.
  ℹtem \textbf{Debt Called In} −− A faction to which the Trust owes a Debt Timer makes a demand that endangers the party.
  ℹtem \textbf{Reputation} −− The Trust gains a reputation (good or bad). High reputation grants +1 die on social rolls within its sphere; low reputation attracts enemies.
\end{itemize}

\paragraph{Example Trust −− The Silkhand League:}
\begin{itemize}
  ℹtem \textbf{Founded}: 10 XP from party pool (Alexander 3, Lyra 2, Sera 3, Mira 2)
  ℹtem \textbf{Seat}: A warehouse in Silkstrand's dyeworks district
  ℹtem \textbf{Initial Assets}: None (spent XP on founding)
  ℹtem \textbf{Debt Timer }: 4 (they owe the local harbor master)
  ℹtem \textbf{Members}: The four PCs
\end{itemize}
After a few sessions, they contribute another 8 XP to the pool and buy a \textbf{Standard Asset}: a spy network among the port's Sidhi pilots. They also add a \textbf{Trust Follower} (Cap 2, a runner named Tema). Now when Lyra needs to know a ship's schedule, she marks 1 Obligation to the Trust and asks Tema. When the Trust's Obligation to her reaches capacity, the Trust will demand that she intercept a rival's cargo −− a mission that becomes the next session's adventure.

\section{Optional: Converting XP to Narrative Wealth}
łabel{sec:xp−narrative−wealth}

Some tables may want a concrete way to spend XP on temporary advantages without acquiring permanent assets.

\begin{center}
\begin{tabular}{ll}
⊤rule
\textbf{XP Spent} & \textbf{Temporary Benefit} \\
∣rule
1 XP & A single use of a Minor Asset (one scene) \\
2 XP & One‑time +2 dice on a roll involving resources (bribe, hire, purchase) \\
3 XP & A Debt Timer segment cleared immediately \\
4 XP & One‑scene use of a Standard Asset (e.g., a ship for a chase) \\
6 XP & One‑mission Follower (Cap 3, must be a willing NPC) \\
⊥tomrule
\end{tabular}
\end{center}

These expenditures represent short‑term investment of effort and contacts −− ``burning'' XP for immediate story gain. They cannot be recovered.

\section{Summary for Players}
łabel{sec:currency−summary}

\begin{itemize}
  ℹtem \textbf{No coins.} Incidental expenses are free.
  ℹtem \textbf{Treasure is leverage.} Minor, Standard, or Major finds that change the story.
  ℹtem \textbf{Significant things cost XP, barter, Strings, or Obligation.}
  ℹtem \textbf{Pool XP} for party assets.
  ℹtem \textbf{Form a Trust} −− a player‑owned terrestrial patron −− for shared power and shared trouble.
  ℹtem \textbf{Use Strings as currency.} Collect favors, trade them, spend them.
  ℹtem \textbf{Debt is story.} Accepting a Debt Timer is not failure; it's future adventure.
\end{itemize}

\begin{tcolorbox}[colback=green!5!white,colframe=green!70!black,title=The Grey Wanderer's Ledger]
\emph{``I have carried silver across a dozen borders and buried it in as many deserts. The only wealth that never tarnishes is a promise witnessed by honest lips. Spend your XP on what you will hold −− a safehouse, a follower, a Trust. The rest is wind.''}
\end{tcolorbox}